% SCP10 arXiv:1001.3807v3, Lemma 6.14, paper_v3.tex:2199--2214.
% Candidate finite-rectangle specialization. Lean compilation is blocked by
% unavailable pinned Mathlib cache; intentionally no leanok or lean tags yet.
\begin{theorem}[Deformation inside an embedded rectangular dual patch]%
    \label{thm:peps_dual_rectangle}
    \notready
    \uses{def:peps_dual_flux_homotopy}
    Fix a dual plaquette $(x_0,y_0)$ and integers $0\leq W<w$,
    $0\leq H<h$. Embed the closed rectangle
    $D=\{0,\ldots,W\}\times\{0,\ldots,H\}$ by
    $\iota(i,j)=(x_0+i,y_0+j)$ modulo $(w,h)$, and put
    $R=\{\iota(i+1,j+1):0\leq i<W,\ 0\leq j<H\}$.
    Every two labelled unit-step paths in $D$ with the same lifted
    endpoints have projections related by $\sim_R$.
    This is the rectangular bulk specialization of
    \cite[Lemma~6.14]{Schuch2010PEPS}.
\end{theorem}
\begin{proof}
    Let $C_{i,j}$ run along the bottom row from $(0,0)$ to $(i,0)$,
    then up to $(i,j)$. Interchanging an east edge with a column of
    height $j$ uses exactly the $j$ adjacent squares. Hence
    $C_a e\sim_R C_b$ for an east edge $e:a\to b$.
    The north-edge identity holds directly; west and south follow
    by cancelling an edge with its reverse. Induction on the steps
    of $p:a\to b$ gives $C_a p\sim_R C_b$.
    Thus $C_a p\sim_R C_a q$; prepend $C_a^{-1}$ and cancel
    its successive backtracks. The bounds $W<w$ and $H<h$ make
    $\iota$ injective on $D$. No periodic-endpoint identification
    replaces equality of the lifted endpoints.
\end{proof}

% Ink-to-index: all dotted arrows are dual geometric routes, with no tensor
% indices or physical legs. In each panel a=(0,0), b=(1,2); vertical drawing
% rows decrease as lifted j increases. The two square centers s_1=(1,1),
% s_2=(1,2) are primal-site labels in the lower-left-plaquette convention.
% The rectangle is two squares high and one square wide. The left route is
% N,N,E; the right route E,N,N. Only these two primal sites are swept.
% This depicts the column/east exchange, not the whole arbitrary-path proof.
\begin{tenkzequation}
    \tenkzkernel
    \begin{tenkz}[size=m, rows={wire,wire,wire,wire,wire,wire,wire}, cols=7, bonds=none]
        \tnwire[kind=string, stroke=dotted, route=orth, dir=to,
            name=rectColumnEast, via={(2,2)}]{(6,2)}{(2,6)}
        \tnmark[form=label, label pos=225]{(6,2)}{$a$}
        \tnmark[form=label, label pos=45]{(2,6)}{$b$}
        \tnmark[form=label]{(5,4)}{$s_1$}
        \tnmark[form=label]{(3,4)}{$s_2$}
    \end{tenkz}
    $\quad\sim_R\quad$
    \begin{tenkz}[size=m, rows={wire,wire,wire,wire,wire,wire,wire}, cols=7, bonds=none]
        \tnwire[kind=string, stroke=dotted, route=orth, dir=to,
            name=rectEastColumn, via={(6,6)}]{(6,2)}{(2,6)}
        \tnmark[form=label, label pos=225]{(6,2)}{$a$}
        \tnmark[form=label, label pos=45]{(2,6)}{$b$}
        \tnmark[form=label]{(5,4)}{$s_1$}
        \tnmark[form=label]{(3,4)}{$s_2$}
    \end{tenkz}
\end{tenkzequation}
